%! AP Statistics — Unit 5, Lesson 5.1 — Lesson Plan
\documentclass[10pt]{article}
\usepackage{apstats-article}
\usepackage{apstats-boxes}

\newcommand{\UnitNumberName}{Unit 5: Sampling Distributions \quad}
\newcommand{\LessonNumberName}{Lesson 5.1: Introducing Statistics --- Why Is My Sample Not Like Yours?}

\begin{document}
\begin{center}
    {\Large\bfseries \CourseName: \SchoolYear \\
    {\small\bfseries \UnitNumberName \LessonNumberName}}
\end{center}

\begin{tcolorbox}[colback=sky, colframe=navy, boxrule=0.9pt, arc=2mm,
    left=2mm, right=2mm, top=1.5mm, bottom=1.5mm]
    \textbf{Primary Objective:} Students will recognize that statistics computed from
    different random samples of the same population will vary, and that this sample-to-sample
    variation is itself a predictable, structured phenomenon --- not simply error
    (Big Idea: VAR).
\end{tcolorbox}

\renewcommand{\arraystretch}{1.6}
\begin{skillbox}[Priority Ideas \& Skills]{goldbox}
\begin{minipage}[t]{0.48\linewidth}
    \textbf{AP Skill 1 --- Selecting Statistical Methods}
    \begin{itemize}
        \item Identify questions suggested by variation in statistics across samples
              from the same population (Skill 1.A).
        \item Recognize that sample statistics vary from sample to sample even when
              drawn from the same population (VAR-1.G, VAR-1.G.1).
        \item Bridge from Unit 4 (probability \& distributions) to Unit 5
              (sampling distributions and inference).
    \end{itemize}
\end{minipage}
\hfill
\begin{minipage}[t]{0.48\linewidth}
    \textbf{Key Understandings}
    \begin{itemize}
        \item Variation in statistics across repeated samples may be random or
              non-random.
        \item The key question for the unit: ``How much should we expect sample
              statistics to vary just by chance?''
        \item Understanding sample-to-sample variability is the foundation for all
              of Units 6--9 (inference).
        \item This lesson is a framing lesson (not assessed in the CED) that sets
              the conceptual stage for the unit.
    \end{itemize}
\end{minipage}
\end{skillbox}

\begin{skillbox}[{Vocabulary, Concepts, \& Theorems}]{greenbox}
\begin{minipage}[t]{0.48\linewidth}
    \begin{itemize}
        \item \textbf{Statistic} --- a number computed from sample data
        \item \textbf{Parameter} --- a fixed (usually unknown) number describing a
              population
        \item \textbf{Sampling variability} --- the natural tendency of a statistic to
              take different values across different samples
    \end{itemize}
\end{minipage}
\hfill
\begin{minipage}[t]{0.48\linewidth}
    \begin{itemize}
        \item \textbf{Sampling distribution} --- the distribution of a statistic over
              all possible samples of a fixed size from a population
        \item \textbf{Random variation} --- fluctuation due to chance alone
        \item \textbf{Bias} --- systematic tendency to over- or under-estimate a
              parameter (preview for Lesson 5.4)
    \end{itemize}
\end{minipage}
\end{skillbox}

\begin{skillbox}[Activate Prior Knowledge \& Spiral Review]{sky}
\begin{minipage}[t]{0.48\linewidth}
    \textbf{Prior Knowledge Check (3--4 min)}
    \begin{itemize}
        \item Review: What is the difference between a \emph{statistic} and a
              \emph{parameter}? (Unit 1)
        \item Connect: Unit 4 built probability models for random variables; Unit 5
              applies that machinery to the statistics we actually compute from data.
        \item Prompt: \textit{``Two students each survey 30 classmates about hours
              of sleep. They get different means. Why?''}
    \end{itemize}
\end{minipage}
\hfill
\begin{minipage}[t]{0.48\linewidth}
    \textbf{Connections Forward}
    \begin{itemize}
        \item Normal distribution as a probability model $\to$ Lesson 5.2
        \item Central Limit Theorem $\to$ Lesson 5.3
        \item Bias and unbiased estimators $\to$ Lesson 5.4
        \item Sampling distributions for proportions $\to$ Lessons 5.5--5.6
        \item Sampling distributions for means $\to$ Lessons 5.7--5.8
        \item Formal inference using these distributions $\to$ Units 6--9
    \end{itemize}
\end{minipage}
\end{skillbox}

\begin{skillbox}[Hook \& Lesson]{sky}
\begin{minipage}[t]{0.48\linewidth}
    \textbf{Hook (5--7 min)}\\
    Display the following scenario:\\[4pt]
    \textit{A polling agency reports that 54\% of a sample of 500 voters support a
    ballot measure. A rival agency samples 500 different voters and gets 51\%. A third
    agency gets 56\%.}\\[4pt]
    Ask: ``Are these agencies lying? Did they make mistakes?''\\[4pt]
    Reveal: All three results can easily come from the same population --- the variation
    is simply \emph{sampling variability}. The real question is: \textbf{How much
    variation should we expect?}
\end{minipage}
\hfill
\begin{minipage}[t]{0.48\linewidth}
    \textbf{Lesson Flow (15 min)}
    \begin{enumerate}
        \item Show three samples of size 10 drawn from a jar of colored chips (or
              simulate via applet). Compute the sample proportion of red for each.
              Record on board.
        \item Ask: \textbf{``Why do these proportions differ?''} (Samples are random;
              different chips get selected.)
        \item Ask: \textbf{``If we took 1000 such samples, what would that collection of
              proportions look like?''} Introduce the term \emph{sampling distribution}.
        \item Key idea (VAR-1.G.1): variation across samples from the same population
              may be random --- and that is \emph{expected}, not a flaw.
        \item Preview: The rest of Unit 5 will describe the exact shape, center, and
              spread of these sampling distributions so we can use them for inference.
    \end{enumerate}
\end{minipage}
\end{skillbox}

\begin{skillbox}[Explicit Instruction \& Active Monitoring]{sky}
\begin{minipage}[t]{0.48\linewidth}
    \textbf{Direct Instruction (8--10 min)}
    \begin{itemize}
        \item Distinguish: distribution of a \emph{population} vs.\ distribution of a
              \emph{sample} vs.\ \emph{sampling distribution} of a statistic.
        \item Use a Desmos applet or class simulation: repeatedly draw samples of size
              $n = 30$ from a known population and plot the sample means. Observe the
              shape of the resulting dot-plot.
        \item Emphasize: the sampling distribution has its own center and spread ---
              these are not random but are determined by the population parameter and $n$.
    \end{itemize}
\end{minipage}
\hfill
\begin{minipage}[t]{0.48\linewidth}
    \textbf{Active Monitoring}
    \begin{itemize}
        \item Listen for students conflating ``the sample'' with ``the sampling
              distribution'' --- redirect: \textit{``One sample gives one dot; the
              sampling distribution is the whole collection of dots.''}
        \item Cold-call: ``Name one question about a population that sampling
              variability makes harder to answer.''
        \item Check warm-up for confusion between statistic and parameter.
    \end{itemize}
\end{minipage}
\end{skillbox}

\begin{skillbox}[Group Work \& Differentiation]{redbox}
\begin{minipage}[t]{0.48\linewidth}
    \textbf{Partner Activity (8--10 min)}\\
    Four scenario cards (each describes a study with two different sample results):
    \begin{enumerate}
        \item State whether the difference is likely due to sampling variability or a
              real population difference.
        \item Identify the statistic and (if stated) the parameter.
        \item Write a question the sampling variability raises that we cannot yet answer.
    \end{enumerate}
\end{minipage}
\hfill
\begin{minipage}[t]{0.48\linewidth}
    \textbf{Differentiation}
    \begin{itemize}
        \item \textit{Support (Tier R)}: Anchor sentence --- ``A statistic changes from
              sample to sample. A parameter stays fixed.'' Provide sentence frames.
        \item \textit{On-grade (Tier A)}: Explain conceptually why increasing sample
              size might reduce sampling variability.
        \item \textit{Extension (Tier E)}: Sketch the sampling distribution of the
              sample mean for $n = 5$ vs.\ $n = 50$. Predict: center? spread? shape?
        \item \textit{ELL}: Bilingual vocabulary card for statistic / parameter /
              sampling variability.
    \end{itemize}
\end{minipage}
\end{skillbox}

\begin{skillbox}[Individual Work \& Assessment]{redbox}
\begin{minipage}[t]{0.48\linewidth}
    \textbf{Exit Ticket (5 min)}
    \begin{enumerate}
        \item A nationwide survey of 1,000 adults finds that 62\% support a new policy.
              A local survey of 50 adults finds 55\%. Explain why these results are not
              necessarily contradictory.
        \item What is the difference between a \emph{statistic} and a \emph{parameter}?
              Give one example of each from the scenario above.
        \item In your own words, what is a sampling distribution?
    \end{enumerate}
\end{minipage}
\hfill
\begin{minipage}[t]{0.48\linewidth}
    \textbf{AP-Style Check}\\
    \textit{A researcher takes repeated random samples of size 40 from a large population
    and records the sample mean each time. The collection of all these sample means is
    best described as:}
    \begin{enumerate}[label=(\Alph*)]
        \item The population distribution
        \item A single sample distribution
        \item The sampling distribution of the sample mean
        \item The parameter of interest
    \end{enumerate}
    Answer: (C)
\end{minipage}
\end{skillbox}

\begin{skillbox}[Reinforcement \& Extension]{goldbox}
\begin{minipage}[t]{0.48\linewidth}
    \textbf{Homework / Practice}
    \begin{itemize}
        \item For three real-world scenarios, identify the statistic, the parameter, and
              whether sampling variability could explain any observed difference between
              two reported statistics.
        \item Reflection: \textit{``Why is it important to understand sampling variability
              before performing formal inference?''}
    \end{itemize}
\end{minipage}
\hfill
\begin{minipage}[t]{0.48\linewidth}
    \textbf{Extension / Preview}
    \begin{itemize}
        \item \textit{Extension}: Find a news article reporting two poll results on the
              same topic. Evaluate whether the difference is likely due to sampling
              variability or a real population difference.
        \item \textit{Preview}: Lesson 5.2 revisits the normal distribution --- the
              key tool for describing what sampling distributions look like.
    \end{itemize}
\end{minipage}
\end{skillbox}

\end{document}
